\documentclass[../RFC-HDFG-2026-001.tex]{subfiles}

\begin{document}

%% ======================================================================
%%  Section 7 – Built-in Filter Parameters
%% ======================================================================
\section{Built-in Filter Parameters}
\label{sec:builtin-filters}

Each built-in filter implements \texttt{set\_config} and \texttt{get\_config} so that the
string API is fully functional without third-party plugins.

For built-in filters that have dedicated \texttt{H5Pset\_*} functions with non-trivial
validation logic (particularly SZIP and Scale-Offset), the \texttt{set\_config}
callback does \emph{not} call the public setter.  Doing so would re-enter the
property-list layer that invoked the callback in the first place, and the public
setters take a \texttt{hid\_t} DCPL rather than the \texttt{cd\_values} array the
callback must fill.  Instead, the validation logic in each such
\texttt{H5Pset\_*} implementation is factored out into a package-internal static
helper (e.g., an \texttt{H5Z\_\_szip\_parms\_check()} extracted from
\texttt{H5Pset\_szip}), and both paths call that one helper.  The
\texttt{set\_config} callback then writes \texttt{cd\_values} itself.  Sharing the
helper---rather than either duplicating the checks or delegating to the public
API---is what prevents divergence between the two configuration paths and
ensures that validation constraints are applied consistently regardless of
whether the filter was configured via \texttt{H5Pset\_szip} or
\texttt{H5Pappend\_filter}.

Throughout this section, ``does not call \texttt{H5Pset\_*}'' should be read in
this sense: the public API is not invoked, but the validation rules it enforces
are, via the shared helper.

\verysubsection{Implementation note}
Each built-in filter is converted to \texttt{H5Z\_class3\_t} in its existing
source file (e.g., \texttt{H5Zdeflate.c}, \texttt{H5Zszip.c}), with
\texttt{set\_config}, \texttt{get\_config}, and a short \texttt{description}
string added to the struct definition.  The \texttt{description} is
informational only (used by introspection tools; see
\SectionRef{sec:class3-struct}) and is set to a one-line summary of the
algorithm (e.g., \texttt{"Deflate (zlib) general-purpose lossless
compression"} for the deflate filter).  This is an entirely internal change;
the existing \texttt{H5Zregister()} path (\SectionRef{sec:class3}) handles
\texttt{H5Z\_class3\_t} registrations without modification.

\subsection{Deflate}

\begin{itemize}
\item \texttt{level} (integer 0--9, default 6): Compression level.
\item Example: \texttt{level = 6}
\end{itemize}

\subsection{Shuffle}

No parameters. \texttt{params} should be \texttt{NULL}.

\subsection{Fletcher32}

No parameters. \texttt{params} should be \texttt{NULL}.

\subsection{SZIP}

\begin{itemize}
\item \texttt{coding} (\texttt{"entropy"} or \texttt{"nn"}, default \texttt{"nn"}): Coding
  method. Value must be a quoted string.
\item \texttt{pixels\_per\_block} (even integer 2--32, default 16): Pixels per block.
\item The \texttt{set\_config} callback (\texttt{H5Z\_\_szip\_set\_config} in
  \texttt{H5Zszip.c}) parses the parameter string directly using
  \texttt{H5Zconfig\_get\_str} and \texttt{H5Zconfig\_get\_int}, validates the
  values, translates \texttt{coding} into the appropriate option mask
  (\texttt{H5\_SZIP\_EC\_OPTION\_MASK} or \texttt{H5\_SZIP\_NN\_OPTION\_MASK}),
  and populates \texttt{cd\_values} directly.  Validation checks that
  \texttt{pixels\_per\_block} is an even integer in \texttt{[2,~32]}, by calling
  the shared internal helper \texttt{H5Z\_\_szip\_parms\_check()} factored out of
  \texttt{H5Pset\_szip} rather than the public \texttt{H5Pset\_szip} API itself.
\item Example: \texttt{coding = "entropy", pixels\_per\_block = 8}
\end{itemize}

\subsection{N-Bit}

No user-facing parameters (derived from datatype). \texttt{params} should be
\texttt{NULL}.

\subsection{Scale-Offset}

\begin{itemize}
\item \texttt{scale\_type} (\texttt{"float\_dscale"}, \texttt{"float\_escale"}, or
  \texttt{"int"}): Scale type. Required. Value must be a quoted string.
\item \texttt{scale\_factor} (non-negative integer): Scale factor. Required.
\item The \texttt{set\_config} callback (\texttt{H5Z\_\_scaleoffset\_set\_config}
  in \texttt{H5Zscaleoffset.c}) parses the parameter string directly using
  \texttt{H5Zconfig\_get\_str} and \texttt{H5Zconfig\_get\_int}, translates
  \texttt{scale\_type} into the corresponding \texttt{H5Z\_SO\_*} enumerator,
  validates that \texttt{scale\_factor} is non-negative---via the shared internal
  helper \texttt{H5Z\_\_scaleoffset\_parms\_check()} factored out of
  \texttt{H5Pset\_scaleoffset} rather than the public
  \texttt{H5Pset\_scaleoffset} API itself---and populates \texttt{cd\_values}
  directly.  The callback returns \texttt{H5E\_BADVALUE} if either key is
  absent.
\item Example: \texttt{scale\_type = "float\_dscale", scale\_factor = 4}
\end{itemize}

\subsection{Example: Blosc2 (Third-Party Filter with Sub-Compressor)}
\label{sec:blosc-example}

Blosc2 (filter ID 32001) is a third-party blocking compressor that wraps an
internal codec---such as \texttt{zlib}, \texttt{lz4}, or \texttt{zstd}---and
applies its own shuffle and blocking logic on top.  It is not a built-in HDF5
filter, but it illustrates the use of nested inline table parameters
(\SectionRef{sec:param-overview}) for filters that accept per-sub-compressor
configuration alongside their own top-level parameters.

The recommended parameter string form uses dotted keys for the sub-compressor:

\begin{minted}{text}
compressor.name = "zlib", compressor.level = 6, shuffle = 1, blocksize = 0
\end{minted}

The inline table form is equivalent:

\begin{minted}{text}
compressor = {name = "zlib", level = 6}, shuffle = 1, blocksize = 0
\end{minted}

\noindent Representative parameters a Blosc2 \texttt{set\_config} might support:

\begin{itemize}
\item \texttt{compressor.name} (string): Internal codec.
  Accepted values: \texttt{"blosclz"}, \texttt{"lz4"}, \texttt{"lz4hc"},
  \texttt{"zlib"}, \texttt{"zstd"}.  Default: \texttt{"blosclz"}.
\item \texttt{compressor.level} (integer 0--9): Compression level passed to
  the internal codec.  Default: \texttt{5}.
\item \texttt{shuffle} (integer 0--2): Shuffle type applied before compression.
  \texttt{0}~=~none, \texttt{1}~=~byte shuffle, \texttt{2}~=~bit shuffle.
  Default: \texttt{1}.
\item \texttt{blocksize} (integer $\geq 0$): Block size in bytes; \texttt{0}
  selects an automatic size.  Default: \texttt{0}.
\item \texttt{nthreads} (integer $\geq 1$): Compression threads.  Default:
  \texttt{1}.
\end{itemize}

\noindent The \texttt{set\_config} callback reads the sub-compressor fields via
dotted-key accessors:

\begin{minted}{c}
char   codec[32];
int64_t clevel, shuffle, blocksize, nthreads;

H5Zconfig_get_str (params, "compressor.name",  codec, sizeof(codec));
H5Zconfig_get_int (params, "compressor.level", &clevel);
H5Zconfig_get_int (params, "shuffle",           &shuffle);
H5Zconfig_get_int (params, "blocksize",         &blocksize);
H5Zconfig_get_int (params, "nthreads",          &nthreads);
\end{minted}

\noindent The same dotted-key form applies whether the caller wrote
\texttt{compressor.name = "zlib"} or \texttt{compressor = \{name = "zlib"\}};
the parser normalises both into the same key path before \texttt{set\_config}
is invoked.


%% ======================================================================
%%  Section 8 – Thread Safety and Concurrency
%% ======================================================================
\section{Thread Safety and Concurrency}
\label{sec:thread-safety}

\subsection{Plugin Loading Protection}

The on-demand plugin scan triggered by \texttt{H5PL\_load} performs filesystem
I/O (directory iteration and shared-library opens). In thread-safe builds of
HDF5 (configured with \texttt{--enable-threadsafe}), this I/O is serialized
under the global API lock.

The existing HDF5 global API lock serializes all public API entry points.
This is sufficient to prevent concurrent plugin loading races from user-facing
calls. No new locking mechanisms are introduced.

\verysubsection{Future-proofing note}
When many threads simultaneously encounter an unknown filter ID at startup,
all threads queue behind the first thread performing the plugin scan, producing
a convoy effect. Reducing lock granularity for the scan itself---for example
by allowing concurrent read-only scans while only serializing the final plugin
table insertion---would benefit throughput if the library moves toward
finer-grained concurrency.

\subsection{Parallel HDF5 (MPI) Considerations}
\label{sec:parallel}

In Parallel HDF5, all MPI ranks participating in collective dataset creation
present identical dataset creation property lists (DCPLs). Non-deterministic
plugin loading across ranks (e.g., due to filesystem latency or differing
\texttt{HDF5\_PLUGIN\_PATH} settings) could cause one rank to fail to load a
required plugin while others succeed.

This is not a new problem introduced by this design---the same class of
issue exists today when any rank fails to load a required plugin, and there
are already several runtime-specific mechanisms (environment variables,
plugin path configuration, dynamic library search order) that can cause
per-rank divergence.

An additional parallel-specific concern is that parameter string parsing is
performed independently on every rank, even though the resulting
\texttt{cd\_values} are identical across ranks and the work need only be done
once.  This is an inherent inefficiency of the current design, though it
should be noted that the existing plugin interface was not designed with
parallel use in mind and this RFC does not claim to address that broader
limitation.

A further concern specific to floating-point parameters is
\textit{strtod} determinism on heterogeneous clusters.  On a cluster mixing
CPU architectures (e.g., x86\_64 and ARM64), the conversion of a decimal
string such as \texttt{"rate=3.55555555"} to a \texttt{double} may produce
different bit patterns if the platform \texttt{strtod} implementations are
not both correctly-rounded per IEEE~754 --- which
\SectionRef{sec:float-precision} makes a normative requirement precisely
because a 1~ulp divergence is not merely a parallel-consistency matter, but
also changes the value a filter receives.  The following applies to
platforms that do not meet that requirement.  The stored configuration string
(\SectionRef{sec:pline-v3}) is identical across ranks by construction, but
the \texttt{cd\_values} produced by each rank's \texttt{set\_config} become
deterministic only after the string-to-\texttt{double} conversion is in
hand, so a per-rank difference in that conversion step propagates into
\texttt{cd\_values} and causes a DCPL mismatch at the collective
\texttt{H5Dcreate} call.  Applications running on heterogeneous clusters
with floating-point filter parameters should obtain a single parsed DCPL on
one rank and broadcast it to all others (e.g., via serialisation through
\texttt{H5Pencode}/\texttt{H5Pdecode}) rather than having each rank
independently invoke \texttt{H5Pappend\_filter}.

\textbf{Note:} \texttt{H5Pdecode} reconstructs a DCPL from the serialised
bytes by copying the stored \texttt{cd\_values} directly; the filter's
\texttt{set\_config} callback is not re-invoked.  The broadcast pattern is
therefore safe on heterogeneous clusters: all ranks obtain an identical
DCPL from the same binary buffer.

\verysubsection{Diagnosing rank-inconsistent DCPLs}
A DCPL mismatch across MPI ranks at \texttt{H5Dcreate} time produces an
\texttt{H5E\_CANTCREATE} error in debug builds but can manifest as a hang
or silent data corruption in optimised builds depending on the MPI
implementation.  Applications can diagnose rank divergence manually:
serialise the DCPL with \texttt{H5Pencode} on every rank, hash the result,
\texttt{MPI\_Allreduce} the hashes, and abort collectively before
\texttt{H5Dcreate} if any hash differs.

\textbf{Future work:} A dedicated environment variable
\texttt{HDF5\_MPI\_CHECK\_DCPL} is proposed for a future release; it would
automate the above check at each collective \texttt{H5Dcreate} and emit a
warning when any rank's hash differs.

\verysubsection{Failure mode and plugin loading order}

\noindent\textbf{Decision:} \texttt{H5Pappend\_filter} loads the plugin
eagerly at call time and invokes \texttt{set\_config} immediately.  Unlike
\texttt{H5Pset\_filter}, which stores raw \texttt{cd\_values} and defers
all plugin interaction, \texttt{H5Pappend\_filter} must run the filter's
\texttt{set\_config} callback to translate the parameter string into
\texttt{cd\_values}; this requires the plugin to be registered.  If the
plugin is unavailable, \texttt{H5Pappend\_filter} returns \texttt{H5E\_NOFILTER}
at the call site rather than deferring failure to \texttt{H5Dcreate} or
I/O time.  For Parallel HDF5, all ranks call \texttt{H5Pappend\_filter}
with identical arguments; a failure on any rank is handled locally before
any collective \texttt{H5Dcreate} is issued.  Applications that want to
probe filter availability before committing to a DCPL configuration should
call \texttt{H5Zfilter\_avail} first.

\subsection{High-Level Language Binding Considerations (h5py, PyTables, etc.)}
\label{sec:hll-bindings}

High-level language bindings such as h5py, PyTables, and netCDF4-Python
call \texttt{H5Pset\_filter} directly with raw \texttt{cd\_values} constructed
on the Python/Cython side.  Because \texttt{H5Pset\_filter} is unchanged by
this RFC, these bindings require \textbf{no changes and no recompilation}.

Should h5py or another binding wish to expose the string form to Python
callers, it would call \texttt{H5Pappend\_filter} from Cython by
declaring \texttt{H5Z\_params\_t} as a Cython \texttt{cdef} struct and
populating it by field assignment (since C99 compound literals are not
available in Cython):

\begin{minted}{python}
# Hypothetical Cython fragment (not part of this RFC)
cdef H5Z_params_t p
p.type  = H5Z_PARAMS_STRING
p.u.str = b"level=6"
H5Pappend_filter(plist_id, H5Z_FILTER_DEFLATE,
                 H5Z_FLAG_MANDATORY, &p)
\end{minted}

\subsection{VOL Connector Interaction}

No changes to the VOL connector interface are required.
\texttt{H5Pappend\_filter} resolves any parameter string into standard
\texttt{cd\_values} before modifying the DCPL, so the execution-relevant
view seen by VOL connectors --- filter ID, flags, \texttt{cd\_values} via
\texttt{H5Pget\_filter2} --- is identical to an entry set via
\texttt{H5Pset\_filter}.  The DCPL additionally retains the parameter
string (\SectionRef{sec:dcpl-retention}); connectors that serialise
property lists via \texttt{H5Pencode}/\texttt{H5Pdecode} carry it
transparently, and connectors that ignore it lose only introspection
fidelity, not correctness.

Thin-client architectures that proxy dataset creation to a remote server
(e.g., REST~VOL) can recover the original parameter string from the DCPL
at any time via \texttt{H5Pget\_filter\_params\_by\_idx}
(\SectionRef{sec:introspection}), since the string is retained rather than
discarded after \texttt{set\_config} resolution.  An earlier draft of this
design, which did not retain the string, required such clients to
negotiate string-form parameters out-of-band; that limitation no longer
applies.


%% ======================================================================
%%  Section 9 – CLI Tool Integration
%% ======================================================================
\section{CLI Tool Integration}
\label{sec:cli-tools}

\subsection{h5repack}

\texttt{h5repack} currently accepts filter specifications via the \texttt{-f} flag. The
existing syntax (\texttt{GZIP}, \texttt{SZIP}, \texttt{SHUF}, \texttt{FLET}, \texttt{NBIT},
\texttt{SOFF}, and \texttt{UD=id,flags,n,v1,...}) continues to work unchanged.

A new form is added that accepts a filter ID and parameter string:

\begin{Verbatim}[fontsize=\scriptsize]
    h5repack -f 'UD=1,0,level=6' input.h5 output.h5
    h5repack -f 'UD=32013,0,mode="rate",rate=3.5' input.h5 output.h5
    h5repack -f '/group/dataset:UD=32013,0,mode="rate",rate=3.5' input.h5 output.h5
\end{Verbatim}

Note: \texttt{mode} is a TOML string value here, so it must be quoted
(\texttt{mode="rate"}); an unquoted \texttt{mode=rate} is not valid TOML and
is rejected by the parser.

The existing \texttt{UD=id,flags,...} syntax is extended to accept a parameter
string in place of the raw \texttt{cd\_values} sequence when the filter is a v3
plugin that implements \texttt{set\_config}. The \texttt{h5repack} parser detects
whether the values following the flags field are a key=value string or a raw
integer sequence and dispatches accordingly---calling \texttt{H5Pappend\_filter}
with \texttt{H5Z\_PARAMS\_STR} for the string form and \texttt{H5Z\_PARAMS\_RAW}
for the raw integer form.
The old built-in aliases (\texttt{GZIP}, \texttt{SHUF}, \texttt{FLET}, \texttt{SOFF})
continue to work unchanged.

\verysubsection{Duplicate filter entries}
HDF5 allows the same filter ID to appear more than once in a pipeline.
When \texttt{h5repack} applies a \texttt{-f} specification that matches a
filter already present in the source pipeline, it appends a new entry
rather than replacing the existing one---this is the same behaviour as the
existing \texttt{UD=} raw-integer form.  Pipeline ordering is fully
preserved: entries are written to the output DCPL in the order they appear
in the source pipeline, followed by any new entries added via \texttt{-f}.
No new behaviour is introduced for this edge case.

\subsection{h5dump}

The \texttt{PARAMS\_STRING} field is displayed only when \texttt{h5dump} is invoked
with \texttt{-p} (the existing flag that enables display of filter and property
details). Output without \texttt{-p} is unchanged. Because \texttt{PARAMS\_STRING} is
gated behind \texttt{-p}, and because format changes of this kind have been
accepted in prior major releases, no additional command-line flag is required.

\verysubsection{Normative output format (with \texttt{-p})}

When \texttt{-p} is active, for each filter in the pipeline, a
\texttt{PARAMS\_STRING} line is appended immediately
after the existing filter line, indented to the same level as other filter
sub-fields:

\begin{verbatim}
    FILTERS {
      COMPRESSION DEFLATE {
        LEVEL 6
        PARAMS_STRING 'level=6'
      }
    }
\end{verbatim}

The value is the literal output of
\texttt{H5Pget\_filter\_params\_by\_idx}, which follows the source
precedence of \SectionRef{sec:introspection}: the configuration string
stored in the version-3 pipeline message when present (exact, and
displayable \emph{without loading the plugin}), otherwise the plugin's
\texttt{get\_config} reconstruction, otherwise the raw \texttt{cd\_values}
formatted as a colon-separated list.  \texttt{PARAMS\_STRING} is therefore
displayed for every filter entry.  The value is enclosed in single
quotes rather than double quotes: the value is TOML text, whose own string
values are double-quoted (e.g. \texttt{mode = "rate"}), so wrapping in
single quotes avoids escaping the common case. A TOML double-quoted string
value may itself contain a literal single quote (e.g. \texttt{"it's"}); if
the parameter string contains one, it is escaped as
\texttt{\textbackslash\textquotesingle} in the h5dump output.

The format change is confined to the \texttt{FILTERS} block; all other
\texttt{h5dump} output is unaffected.

\verysubsection{Hexadecimal annotation of binary-exact floats}
\label{sec:h5dump-hex-annotation}

Canonicalisation preserves a float's \emph{value} exactly but not its text
(\SectionRef{sec:config-canonical-form}): a caller who writes
\texttt{accuracy = 0x1p-36} reads back
\texttt{accuracy = 1.4551915228366852e-11}.  The decimal is unrecognisable
as a power of two, which is the form a user tuning a
quantising filter reasons in.  The original text cannot be carried in the
stored bytes --- a conforming TOML parser discards comments, so a
comment-aware reader and a stock one would derive different values from
identical bytes --- so the rendering is produced at display time from the
parsed \texttt{double}, and nothing is added to the file.

When \texttt{-p} is active, \texttt{h5dump} appends a trailing comment
naming each float in the parameter string whose value is a small multiple
of a power of two:

\begin{verbatim}
    FILTERS {
      USER_DEFINED_FILTER {
        FILTER_ID 32013
        COMMENT zfp
        PARAMS_STRING 'mode = "accuracy", accuracy = 1.4551915228366852e-11'
          # accuracy = 0x1p-36
      }
    }
\end{verbatim}

\begin{notebox}
\textbf{Normative rules:}
\begin{itemize}
  \item \textbf{Placement.}  The comment is emitted \emph{outside} the
    single-quoted value, never within it.  The quoted text remains the
    literal output of \texttt{H5Pget\_filter\_params\_by\_idx}, byte for
    byte, so a script that extracts the quoted span still recovers exactly
    the stored bytes.
  \item \textbf{Selection.}  A float is annotated only when its significand
    fits in 13 bits --- equivalently, when the value is $m \cdot 2^{e}$ for
    an integer $m$ with $|m| < 2^{13}$, which is at most three hexadecimal
    fraction digits.  Annotating merely because the hex form is
    \emph{shorter} than the decimal is not the rule and would be nearly
    vacuous: \texttt{0.1} renders as \texttt{0x1.999999999999ap-4}, 20
    characters against the canonical form's 22, so almost every float would
    acquire a comment that tells the reader nothing.  The 13-bit bound
    admits the values this exists for --- exact powers of two, and simple
    multiples such as \texttt{0x1.8p-36} --- and excludes values whose hex
    form is merely a second unreadable spelling.
  \item \textbf{Rendering.}  \texttt{h5dump} formats the hexadecimal text
    itself; it does \emph{not} print \texttt{\%a} output verbatim.  C17
    \S7.21.6.1p8 requires only that a precision-free \texttt{\%a} be
    ``sufficient for an exact representation'' when \texttt{FLT\_RADIX} is a
    power of 2, and constrains the leading digit only to be nonzero for
    normalized values, so \texttt{0x1.8p-36},
    \texttt{0x1.8000000000000p-36}, and \texttt{0x3p-37} are all conforming
    renderings of one \texttt{double} --- glibc trims trailing zeros and the
    Microsoft CRT does not.  \texttt{h5dump} output is compared byte for
    byte against expected files, so the tool emits a canonical spelling: a
    normalized leading digit, trailing zeros of the fraction removed, the
    fraction and its decimal point omitted entirely when zero, lowercase
    \texttt{0x} and \texttt{p}, and a signed exponent.
  \item \textbf{Already-hexadecimal values are not annotated.}  A literal
    the reader can already see in hexadecimal needs no second spelling.
    This does not arise for a canonical stored string, which by
    construction has no hex-float tokens left
    (\SectionRef{sec:config-canonical-form}), but it does for a
    \texttt{get\_config} reconstruction, which is not canonicalised.
  \item \textbf{Scope.}  Display only.  The annotation never affects the
    stored bytes, \texttt{H5Pget\_filter\_params\_by\_idx}, or any value
    passed to \texttt{set\_config}, and its absence is never diagnostic ---
    a float without a comment is simply not a small multiple of a power of
    two.  The rule applies wherever a tool displays a parameter string;
    \texttt{h5ls~-v} prints raw \texttt{cd\_values} today and would adopt it
    along with \texttt{PARAMS\_STRING} display, which this RFC does not
    propose for that tool.
  \item \textbf{Mechanism.}  The tool locates the floats by a lexical scan
    of the parameter string, skipping quoted strings and comments exactly as
    the hex-float rewrite does, and taking the preceding bare key --- dotted
    keys included --- as the name to echo.  No key-enumeration API is
    required, and none is proposed: the typed accessors
    (\texttt{H5Zconfig\_get\_double} and its siblings) are lookup-by-name,
    and \texttt{h5dump} does not know the key names of a filter it has not
    loaded.  A scan is admissible precisely because the annotation is
    display-only --- it cannot alter a value, only fail to comment on one.
\end{itemize}
\end{notebox}

Generating the hexadecimal text rather than delegating to \texttt{\%a} also
improves the subnormal case: \texttt{frexp} normalizes, so
$2^{-1074}$ annotates as \texttt{0x1p-1074} rather than the
\texttt{0x0.0000000000001p-1022} that a denormalized \texttt{\%a} rendering
produces.  Both are exact; only the first is legible, and both re-read to the
identical \texttt{double}.

The reciprocal direction is the one that carries the workflow: hex-float
literals are accepted on \emph{input} everywhere a parameter string is
(\SectionRef{sec:float-precision}), including \texttt{h5repack~-f}, so a
sweep can be driven directly from the shell without converting tolerances
to decimal first:

\begin{verbatim}
    for exp in $(seq -1 -1 -64); do
      h5repack -f "UD=32013,0,mode = \"accuracy\", accuracy = 0x1p$exp" \
               in.h5 out_$exp.h5
    done
\end{verbatim}

\verysubsection{DESCRIPTION}

When \texttt{-p} is active, a \texttt{DESCRIPTION} line is appended after
\texttt{PARAMS\_STRING}, nested inside the entry of the filter it describes,
for filters whose class registration supplies a non-NULL \texttt{description}
(via \texttt{H5Zget\_filter\_class\_info}):

\begin{verbatim}
    FILTERS {
      USER_DEFINED_FILTER {
        FILTER_ID 32013
        COMMENT zfp
        PARAMS { ... }
        PARAMS_STRING 'mode="rate",rate=3.5'
        DESCRIPTION "ZFP compression"
      }
    }
\end{verbatim}

Unlike \texttt{COMMENT} (the filter's canonical name, which is persisted in
the pipeline message and always displays) and \texttt{PARAMS\_STRING} (which
falls back to a raw \texttt{cd\_values} listing when \texttt{get\_config} is
unavailable), \texttt{DESCRIPTION} has no on-disk fallback: it is a
best-effort, registry-sourced label read from the filter class's in-memory
registration at the time \texttt{h5dump} runs. If the filter plugin is not
registered or loadable on the machine running \texttt{h5dump}, the
\texttt{DESCRIPTION} line is omitted entirely. Running \texttt{h5dump -p} on
the same file on two machines -- one with the plugin installed, one without
-- can therefore produce different output for this line. Applications and
scripts that parse \texttt{h5dump -p} output must treat \texttt{DESCRIPTION}
as optional and machine-dependent, and should rely on \texttt{COMMENT}
instead for a stable, plugin-independent label.

%% ======================================================================
%%  Section 10 – Language Bindings
%% ======================================================================
\section{Language Bindings}
\label{sec:bindings}

\subsection{Fortran}
\label{sec:fortran-bindings}

Fortran's ISO\_C\_BINDING interop cannot represent a C union, so
\texttt{H5Z\_params\_t} cannot be expressed as a \texttt{BIND(C)} derived
type.  Two thin C helper functions are therefore added to
\texttt{H5Zf.c} --- one per calling mode --- each of which constructs the
appropriate \texttt{H5Z\_params\_t} value and calls
\texttt{H5Pappend\_filter}:

\begin{minted}{c}
/* In H5Zf.c -- C-side entry points for Fortran BIND(C).
   H5Z_params_t contains a union; Fortran cannot represent unions, so
   these helpers construct the struct by field assignment rather than
   using the H5Z_PARAMS_STR / H5Z_PARAMS_RAW compound-literal macros. */

herr_t H5Pappend_filter_str_c(hid_t plist, H5Z_filter_t id,
                               unsigned flags, const char *params)
{
    H5Z_params_t p;
    p.type  = H5Z_PARAMS_STRING;
    p.u.str = params;
    return H5Pappend_filter(plist, id, flags, &p);
}

herr_t H5Pappend_filter_raw_c(hid_t plist, H5Z_filter_t id,
                               unsigned flags, size_t cd_nelmts,
                               const unsigned *cd_values)
{
    H5Z_params_t p;
    p.type            = H5Z_PARAMS_CDVALUES;
    p.u.raw.cd_nelmts = cd_nelmts;
    p.u.raw.cd_values = cd_values;
    return H5Pappend_filter(plist, id, flags, &p);
}
\end{minted}

The Fortran binding exposes a generic interface
\texttt{h5pappend\_filter\_f} with two specific procedures.  Following the
established HDF5 Fortran convention, each specific is a standard wrapper
subroutine containing an internal \texttt{INTERFACE} block that binds to the
corresponding C helper.  The wrapper handles string conversion (NUL
termination) before forwarding to C.

\begin{minted}{fortran}
! In H5Pff.F90 -- generic interface
INTERFACE h5pappend_filter_f
  MODULE PROCEDURE h5pappend_filter_str_f   ! string form
  MODULE PROCEDURE h5pappend_filter_raw_f   ! cd_values form
END INTERFACE

! String form -- wrapper calling H5Pappend_filter_str_c via internal INTERFACE
SUBROUTINE h5pappend_filter_str_f(prp_id, filter, flags, params, hdferr)
  IMPLICIT NONE
  INTEGER(HID_T),   INTENT(IN)  :: prp_id
  INTEGER,          INTENT(IN)  :: filter
  INTEGER,          INTENT(IN)  :: flags
  CHARACTER(LEN=*), INTENT(IN)  :: params
  INTEGER,          INTENT(OUT) :: hdferr

  CHARACTER(LEN=LEN_TRIM(params)+1, KIND=C_CHAR) :: c_params
  INTERFACE
    INTEGER(C_INT) FUNCTION H5Pappend_filter_str_c( &
        plist_id, filter_c, flags_c, params_c) &
        BIND(C, NAME='H5Pappend_filter_str_c')
      IMPORT :: HID_T, C_INT, C_CHAR
      INTEGER(HID_T), VALUE                            :: plist_id
      INTEGER(C_INT), VALUE                            :: filter_c  ! H5Z_filter_t
      INTEGER(C_INT), VALUE                            :: flags_c   ! unsigned
      CHARACTER(KIND=C_CHAR), DIMENSION(*), INTENT(IN) :: params_c
    END FUNCTION H5Pappend_filter_str_c
  END INTERFACE
  c_params = TRIM(params)//C_NULL_CHAR
  hdferr   = INT(H5Pappend_filter_str_c(prp_id, INT(filter,C_INT), &
                                        INT(flags,C_INT), c_params))
END SUBROUTINE h5pappend_filter_str_f

! Raw cd_values form -- wrapper calling H5Pappend_filter_raw_c
SUBROUTINE h5pappend_filter_raw_f(prp_id, filter, flags, cd_nelmts, cd_values, hdferr)
  IMPLICIT NONE
  INTEGER(HID_T),  INTENT(IN)               :: prp_id
  INTEGER,         INTENT(IN)               :: filter
  INTEGER,         INTENT(IN)               :: flags
  INTEGER(SIZE_T), INTENT(IN)               :: cd_nelmts
  INTEGER,         DIMENSION(*), INTENT(IN) :: cd_values
  INTEGER,         INTENT(OUT)              :: hdferr

  ! Default-kind INTEGER is not necessarily C_INT (e.g. under -i8), so the
  ! array is converted element-wise rather than passed through.
  INTEGER(C_INT), DIMENSION(:), ALLOCATABLE :: cd_values_c
  INTERFACE
    INTEGER(C_INT) FUNCTION H5Pappend_filter_raw_c( &
        plist_id, filter_c, flags_c, cd_nelmts_c, cd_vals) &
        BIND(C, NAME='H5Pappend_filter_raw_c')
      IMPORT :: HID_T, C_INT, SIZE_T
      INTEGER(HID_T), VALUE                    :: plist_id
      INTEGER(C_INT), VALUE                    :: filter_c  ! H5Z_filter_t
      INTEGER(C_INT), VALUE                    :: flags_c   ! unsigned
      INTEGER(SIZE_T), VALUE                   :: cd_nelmts_c
      INTEGER(C_INT), DIMENSION(*), INTENT(IN) :: cd_vals
    END FUNCTION H5Pappend_filter_raw_c
  END INTERFACE

  ALLOCATE(cd_values_c(MAX(cd_nelmts,1_SIZE_T)), STAT=hdferr)
  IF (hdferr /= 0) THEN
     hdferr = -1
     RETURN
  END IF
  cd_values_c(1:cd_nelmts) = INT(cd_values(1:cd_nelmts), C_INT)

  hdferr = INT(H5Pappend_filter_raw_c(prp_id, INT(filter,C_INT), &
               INT(flags,C_INT), cd_nelmts, cd_values_c))
  DEALLOCATE(cd_values_c)
END SUBROUTINE h5pappend_filter_raw_f
\end{minted}

\verysubsection{Fortran introspection: \texttt{h5pget\_filter\_params\_by\_idx\_f}}

\texttt{H5Pget\_filter\_params\_by\_idx} (\SectionRef{sec:introspection}) takes
no union argument, so its wrapper binds directly to the public C entry point
and needs no \texttt{H5Zf.c} helper.  The wrapper hides the C size-query
protocol: it calls the C function twice --- once with a null buffer to learn
the length, once to populate an exactly sized temporary --- and copies the
result into the caller's \texttt{CHARACTER} variable.  \texttt{params\_len}
returns the \emph{full} length of the parameter string even when
\texttt{params} was too short to hold it, so a caller can detect truncation by
comparing it against \texttt{LEN(params)}.

\begin{minted}{fortran}
! In H5Pff.F90
SUBROUTINE h5pget_filter_params_by_idx_f(prp_id, filter_idx, params, hdferr, &
                                        params_len)
  IMPLICIT NONE
  INTEGER(HID_T),   INTENT(IN)            :: prp_id
  INTEGER,          INTENT(IN)            :: filter_idx
  CHARACTER(LEN=*), INTENT(OUT)           :: params
  INTEGER,          INTENT(OUT)           :: hdferr
  INTEGER(SIZE_T),  INTENT(OUT), OPTIONAL :: params_len

  CHARACTER(LEN=1,KIND=C_CHAR), DIMENSION(:), ALLOCATABLE, TARGET :: buf
  INTEGER(SIZE_T) :: needed, ncopy, i
  INTEGER(C_INT)  :: rc
  INTERFACE
    INTEGER(C_INT) FUNCTION H5Pget_filter_params_by_idx_c( &
        plist_id, filter_idx_c, params_buf, params_buf_size, params_len_c) &
        BIND(C, NAME='H5Pget_filter_params_by_idx')
      IMPORT :: HID_T, C_INT, C_PTR, SIZE_T
      INTEGER(HID_T),  VALUE :: plist_id
      INTEGER(C_INT),  VALUE :: filter_idx_c     ! unsigned
      TYPE(C_PTR),     VALUE :: params_buf       ! C_NULL_PTR => size query
      INTEGER(SIZE_T), VALUE :: params_buf_size
      INTEGER(SIZE_T)        :: params_len_c     ! by reference
    END FUNCTION H5Pget_filter_params_by_idx_c
  END INTERFACE

  params = ''
  IF (PRESENT(params_len)) params_len = 0_SIZE_T

  ! Pass 1: size query (params_len is set even with a null buffer)
  rc = H5Pget_filter_params_by_idx_c(prp_id, INT(filter_idx,C_INT), &
                                     C_NULL_PTR, 0_SIZE_T, needed)
  IF (rc < 0) THEN
     hdferr = INT(rc)
     RETURN
  END IF
  IF (PRESENT(params_len)) params_len = needed

  ! Pass 2: populate an exactly sized buffer, then copy out
  ALLOCATE(buf(needed+1), STAT=hdferr)
  IF (hdferr /= 0) THEN
     hdferr = -1
     RETURN
  END IF
  rc = H5Pget_filter_params_by_idx_c(prp_id, INT(filter_idx,C_INT), &
                                     C_LOC(buf(1)), needed+1, needed)
  IF (rc >= 0) THEN
     ncopy = MIN(needed, INT(LEN(params), SIZE_T))
     DO i = 1, ncopy
        params(i:i) = buf(i)
     END DO
  END IF
  DEALLOCATE(buf)
  hdferr = INT(rc)
END SUBROUTINE h5pget_filter_params_by_idx_f
\end{minted}

\noindent The returned \texttt{params} is blank-padded to the declared length
of the caller's variable, following the usual Fortran convention; the
NUL-terminated C form is not exposed.  A filter configured through the raw
\texttt{cd\_values} path yields the \texttt{get\_config} reconstruction or the
\texttt{cd\_values} fallback string, exactly as in C
(\SectionRef{sec:introspection}).

\verysubsection{Fortran interoperable type mapping}
\texttt{H5Z\_filter\_t} is \texttt{typedef int} in \texttt{H5Zpublic.h}, so
\texttt{INTEGER(C\_INT)} is its exact interoperable counterpart on every
platform HDF5 supports; no separate \texttt{H5Z\_FILTER\_T} Fortran kind
parameter is needed.  The \texttt{flags} argument is \texttt{unsigned} in C and
is likewise passed as \texttt{INTEGER(C\_INT)}: Fortran has no unsigned integer
type, and this is the convention the existing HDF5 Fortran binding already uses
for \texttt{H5Pset\_filter}'s \texttt{flags}.  The filter flag bits in use
occupy the low bits of the word, so the signed/unsigned reinterpretation is
value-preserving in practice; the C helpers in \texttt{H5Zf.c} cast back to
\texttt{unsigned} on entry.

\textbf{Default-kind \texttt{INTEGER} arguments are converted, never passed
through.}  The user-facing arguments \texttt{filter}, \texttt{flags},
\texttt{filter\_idx}, and \texttt{cd\_values} are declared as default-kind
\texttt{INTEGER} so that callers need not name a specific kind, but default
\texttt{INTEGER} is not guaranteed to be \texttt{C\_INT}: a build using
\texttt{-fdefault-integer-8} or \texttt{-i8} makes it 64-bit while the C
interface still expects 32-bit.  Scalars are therefore converted at the call
site with \texttt{INT(x, C\_INT)}, and the \texttt{cd\_values} array is
converted element-wise into an \texttt{INTEGER(C\_INT)} temporary rather than
passed through --- an unadorned \texttt{INTEGER, DIMENSION(*)} actual argument
would otherwise reach C as an array of the wrong element width, misreading
every element after the first.  When default \texttt{INTEGER} already is
\texttt{C\_INT} the conversion is a plain copy.  The temporary is
\texttt{ALLOCATABLE} rather than an automatic array because \texttt{cd\_nelmts}
is bounded only by \texttt{UINT16\_MAX}
(\SectionRef{sec:callback-specs}); typical filters use fewer than a dozen
elements, so the allocation is negligible against the surrounding property-list
work.

\textbf{Range caveat.}  \texttt{cd\_values} elements are \texttt{unsigned int}
in C, and Fortran cannot represent values above \texttt{HUGE(0)} in a
default-kind \texttt{INTEGER}.  Callers that need the full unsigned range ---
for example a filter whose \texttt{cd\_values} carry the bit pattern of a
\texttt{double} (\SectionRef{sec:cd-packing-convention}) --- pass those
elements as negative two's-complement values, which reach C with the intended
bit pattern.  This is pre-existing behaviour of
\texttt{h5pset\_filter\_f} and is unchanged here; the string form
(\texttt{h5pappend\_filter\_str\_f}) avoids the question entirely, which is
the recommended path for Fortran callers.

\verysubsection{Fortran string passing convention}
Fortran \texttt{CHARACTER(LEN=*)} strings are not NUL-terminated; the
compiler right-pads them with spaces to their declared length.  The C
helper \texttt{H5Pappend\_filter\_str\_c} expects a NUL-terminated
\texttt{const char~*}.  The wrapper subroutine is responsible for trimming
and appending \texttt{C\_NULL\_CHAR} before calling the C helper, following
the same convention used throughout the existing HDF5 Fortran binding.
The wrapper constructs a trimmed, NUL-terminated local variable:

\begin{minted}{fortran}
c_params = TRIM(params)//C_NULL_CHAR
\end{minted}

This correctly handles \texttt{CHARACTER} variables that are wider than the
string literal assigned to them (Fortran right-pads with spaces).  No
additional stripping is required on the C side.

\verysubsection{Fortran usage example}

\begin{minted}{fortran}
USE HDF5
USE ISO_C_BINDING
IMPLICIT NONE
INTEGER(HID_T)        :: plist
INTEGER               :: hdferr
INTEGER, DIMENSION(1) :: cd_vals = [6]
CHARACTER(LEN=256)    :: params
INTEGER(SIZE_T)       :: params_len

CALL h5pcreate_f(H5P_DATASET_CREATE_F, plist, hdferr)

! String form (generic resolves to h5pappend_filter_str_f)
CALL h5pappend_filter_f(plist, H5Z_FILTER_DEFLATE_F, &
                             H5Z_FLAG_MANDATORY_F,        &
                             "level=6", hdferr)

! Raw cd_values form (generic resolves to h5pappend_filter_raw_f)
CALL h5pappend_filter_f(plist, H5Z_FILTER_DEFLATE_F, &
                             H5Z_FLAG_MANDATORY_F,        &
                             1_SIZE_T, cd_vals, hdferr)

! Introspection: read back entry 0's configuration string
CALL h5pget_filter_params_by_idx_f(plist, 0, params, hdferr, params_len)
IF (hdferr >= 0 .AND. params_len <= LEN(params)) &
   WRITE(*,'(A)') "filter 0: "//TRIM(params)     ! -> level = 6
\end{minted}


\subsection{Java}
\label{sec:java-bindings}

Java does not support C unions, and method overloading is the idiomatic
Java mechanism for variant inputs.  Two overloaded static methods are
added to the existing \texttt{HDF5Library} class:

\begin{minted}{java}
// String form
public static native int H5Pappend_filter(
    long   plist_id,
    int    filter_id,
    int    flags,
    String params       // null means "no parameters"
) throws HDF5LibraryException, NullPointerException;

// Raw cd_values form (equivalent to existing H5Pset_filter)
public static native int H5Pappend_filter(
    long   plist_id,
    int    filter_id,
    int    flags,
    int[]  cd_values    // null or zero-length means "no parameters"
) throws HDF5LibraryException;
\end{minted}

Each JNI implementation (in \texttt{H5P.c} on the native side) constructs
the appropriate \texttt{H5Z\_params\_t} value using \texttt{H5Z\_PARAMS\_STR}
or \texttt{H5Z\_PARAMS\_RAW} and calls \texttt{H5Pappend\_filter}.
The \texttt{String} overload converts the Java string to a NUL-terminated
\texttt{const char~*} via \texttt{GetStringUTFChars} before constructing
the \texttt{H5Z\_params\_t}.

The following additional method is added for pipeline introspection:

\begin{minted}{java}
// Wraps H5Pget_filter_params_by_idx
public static native String H5Pget_filter_params_by_idx(
    long plist_id,
    int  filter_idx
) throws HDF5LibraryException;
\end{minted}

\verysubsection{Java usage example}

\begin{minted}{java}
import hdf.hdf5lib.H5;
import hdf.hdf5lib.HDF5Constants;

long plist = H5.H5Pcreate(HDF5Constants.H5P_DATASET_CREATE);

// String form
H5.H5Pappend_filter(plist, HDF5Constants.H5Z_FILTER_DEFLATE,
                         HDF5Constants.H5Z_FLAG_MANDATORY, "level=6");

// Raw cd_values form
H5.H5Pappend_filter(plist, HDF5Constants.H5Z_FILTER_DEFLATE,
                         HDF5Constants.H5Z_FLAG_MANDATORY,
                         new int[]{ 6 });
\end{minted}


\subsection{\texttt{H5Pmodify\_filter\_by\_idx} in the Bindings}
\label{sec:modify-bindings}

\texttt{H5Pmodify\_filter\_by\_idx} (\SectionRef{sec:modify-filter}) takes
the same \texttt{H5Z\_params\_t} as \texttt{H5Pappend\_filter} and needs no
new binding machinery.  Fortran gains
\texttt{h5pmodify\_filter\_by\_idx\_f} as a generic interface over the same
two specific procedures --- one per calling mode --- backed by two more
thin \texttt{H5Zf.c} helpers built on the identical field-by-field struct
construction shown above.  Java gains two overloads of
\texttt{H5Pmodify\_filter\_by\_idx}, taking a \texttt{String} and an
\texttt{int[]} respectively, mirroring the \texttt{H5Pappend\_filter}
overloads.  In both languages the additional argument relative to append is
the \texttt{filter\_idx}, and in both the raw form clears any stored
configuration string exactly as the C API does.

\end{document}
